\documentclass[11pt]{article}
\usepackage{preamble}
\usepackage{tikz}
\usetikzlibrary{arrows.meta,positioning,fit,backgrounds,calc}
\usepackage{titling}
\setlength{\droptitle}{-2em}
\posttitle{\par\end{center}\vspace{-1em}}

\title{Primer: AIRL with Unobserved Heterogeneity (AIRL-Het)\\[4pt]
\large Lee, Sudhir and Wang (2026)}
\author{}
\date{}

\begin{document}
\maketitle
\vspace{-1.5em}

\section*{Setup}

\paragraph{Actions.} $\mathcal{A} = \{0, 1, 2\}$ corresponding to pay (immediate access), wait (delayed free access), and exit.

\paragraph{State.} $s_{ijt} = (C_{ijt},\; g_{ijt},\; P^0_{ijt},\; P^1_{ijt},\; w_j,\; \texttt{pay}_{ij(t-1)},\; X_{ijt})$ where $C$ is a 32-dim content embedding, $g$ is progress through the series, $P^0, P^1$ are price indicators, $w_j$ is the series-level wait-time, $\texttt{pay}_{ij(t-1)}$ indicates whether the previous episode was purchased, and $X$ contains narrative features.

\paragraph{Heterogeneity.} Each user $i$ belongs to a latent segment $z \in \{1, \ldots, K\}$ with $K = 4$. Segment membership is shared across all series consumed by the same user.

\paragraph{Objective.} Recover segment-specific reward functions $r^*(s, a, z)$ and policies $\pi^*(a \mid s, z)$ that rationalize observed trajectories under the maximum-entropy framework:
\[
\max_\pi \; \mathbb{E}\!\left[\sum_{t=0}^{\infty} \gamma^t \big(r(s_t, a_t, z) + \mathcal{H}(\pi(\cdot \mid s_t, z))\big)\right], \quad \gamma = 0.99.
\]

\vspace{-0.5em}
\paragraph{Termination.} Exit ($a = 2$) or series completion ($g = 1$). Both lead to absorbing state $s_{\text{abs}}$ with $r(s_{\text{abs}}, a, z) = 0$ and $V(s_{\text{abs}}) = 0$.

\paragraph{Networks.}\leavevmode
\vspace{-0.5em}

\begin{table}[h]
\centering\small
\begin{tabular}{lll}
\toprule
Component & Parameterizes & Architecture \\
\midrule
$g_\theta(s, a, z)$ & Immediate reward & 2-layer MLP (128, 64), Leaky ReLU \\
$h_\phi(s, z)$ & Shaping potential $\approx V(s)$ & 2-layer MLP (128, 64), Leaky ReLU \\
$\pi_\omega(a \mid s, z)$ & Policy (CCPs) & 2-layer MLP (128, 64), Leaky ReLU \\
$q_\psi(z \mid \tau)$ & Posterior (segment assignment) & Bi-LSTM (128) + MLP (256, 64) \\
\bottomrule
\end{tabular}
\end{table}

\vspace{-0.5em}
\paragraph{Discriminator structure.} The discriminator embeds a Bellman-consistent decomposition:
\begin{align}
f_{\theta,\phi}(s, a, s', z) &= g_\theta(s, a, z) + \gamma\, h_\phi(s', z) - h_\phi(s, z), \label{eq:f}\\[2pt]
D_{\theta,\phi}(s, a, s', z) &= \frac{\exp\!\big(f_{\theta,\phi}(s, a, s', z)\big)}{\exp\!\big(f_{\theta,\phi}(s, a, s', z)\big) + \pi_\omega(a \mid s, z)}. \label{eq:D}
\end{align}
Note: $\phi$ parameterizes the shaping potential $h$, not the $Q$-function. The policy $\pi_\omega$ is a separate network. The $Q$-function is never directly parameterized.

\paragraph{Additional notation.} We suppress subscripts $ijt$ in the algorithm for clarity.
\vspace{-0.5em}

\begin{table}[h!]
\centering\small
\begin{tabular}{ll}
\toprule
Symbol & Meaning \\
\midrule
$\tau$ & Trajectory $(s_0, a_0, s_1, a_1, \ldots)$ for one user in one series \\
$\mathcal{D}_{\text{obs}}$ & Set of observed transitions $(s, a, s')$ from platform data \\
$\mathcal{D}_{\text{sim}}$ & Set of simulated transitions from current policy $\pi_\omega$ \\
$f$ & Discriminator logit (real-valued score) \\
$D$ & Discriminator output $\in [0, 1]$ (probability transition is observed) \\
$\alpha$ & Learning rate for gradient descent \\
$\varepsilon$ & Convergence tolerance for stopping criterion \\
$\lambda_{\text{CR}}$ & Weight on within-user consistency regularization (S3) \\
$\lambda_{\text{MI}}$ & Weight on mutual information reward (S4); set to 0 after posterior converges \\
\bottomrule
\end{tabular}
\end{table}

\paragraph{What each step does.}
\textbf{S1} generates synthetic reading behavior by running the current policy forward through actual book content.
\textbf{S2} is the core step: it updates the reward and value networks so the discriminator can tell real reading behavior from fake. This is where the structural parameters $\theta$ (reward) and $\phi$ (value) are learned.
\textbf{S3} assigns each user to a latent segment, enforcing that the same user behaves consistently across different series.
\textbf{S4} converts the discriminator's output into a scalar reward signal for each simulated transition.
\textbf{S5} updates the policy to maximize these rewards using PPO. This is the forward RL step that would be soft value iteration in a tabular MDP but requires PPO because the state space is high-dimensional.

\newpage
\section*{Algorithm: LSW AIRL with Heterogeneity}

\begin{center}
\begin{tikzpicture}[
    node distance=0.6cm,
    >={Stealth[length=3mm]},
    block/.style={
        rectangle, draw, rounded corners=3pt,
        text width=13.5cm, align=left,
        inner sep=8pt, font=\small
    },
    stepblock/.style={
        rectangle, draw, rounded corners=3pt, fill=gray!8,
        text width=13.5cm, align=left,
        inner sep=8pt, font=\small
    },
    label/.style={
        font=\small\bfseries, anchor=west
    },
    arrow/.style={->, thick, draw=black!70},
    looparrow/.style={->, very thick, draw=blue!60},
]

% Init
\node[block] (init) {
    \textbf{Initialize} parameters $\theta_0, \phi_0, \omega_0, \psi_0$
};

% S1
\node[stepblock, below=0.5cm of init] (s1) {
    \textbf{S1. Simulate Trajectories} \hfill \textit{(generates $\mathcal{D}_{\text{sim}}$)}\\[3pt]
    \textbf{for} $z = 1, \ldots, K$ \textbf{do}\\
    \quad Sample series $j$, set initial state $s_0$\\
    \quad Roll out $\pi_\omega(a \mid s, z)$ episode-by-episode to generate $\tau^{\text{sim}}$\\
    \quad Stop on exit ($a\!=\!2$) or series completion ($g\!=\!1$) $\to$ enter $s_{\text{abs}}$\\
    \textbf{end for}
};

% S2
\node[stepblock, below=0.5cm of s1] (s2) {
    \textbf{S2. Discriminator Update} \hfill \textit{(updates $\theta, \phi$)}\\[3pt]
    \textbf{for} each $(s, a, s')$ in $\mathcal{D}_{\text{obs}} \cup \mathcal{D}_{\text{sim}}$ \textbf{do}\\
    \quad $f \leftarrow g_\theta(s, a, z) + \gamma\, h_\phi(s', z) - h_\phi(s, z)$\\
    \quad $D \leftarrow \exp(f)\,/\,\bigl(\exp(f) + \pi_\omega(a \mid s, z)\bigr)$\\
    \textbf{end for}\\
    $\mathcal{L}_{\text{disc}} \leftarrow -\frac{1}{|\mathcal{D}_{\text{obs}}|}\sum_{\mathcal{D}_{\text{obs}}} \log D \;-\; \frac{1}{|\mathcal{D}_{\text{sim}}|}\sum_{\mathcal{D}_{\text{sim}}} \log(1 - D)$\\
    $\theta, \phi \leftarrow \theta, \phi - \alpha\,\nabla_{\theta,\phi}\,\mathcal{L}_{\text{disc}}$
};

% S3
\node[stepblock, below=0.5cm of s2] (s3) {
    \textbf{S3. Posterior Update} \hfill \textit{(updates $\psi$)}\\[3pt]
    \textbf{for} each trajectory $\tau_i$ \textbf{do}\\
    \quad Compute $q_\psi(z \mid \tau_i)$\\
    \textbf{end for}\\
    $\mathcal{L}_{\text{post}} \leftarrow \mathcal{L}_{\text{classify}} + \lambda_{\text{CR}} \cdot \mathcal{L}_{\text{consistency}}$
    \quad {\footnotesize (same user across series $\to$ similar $z$)}\\
    $\psi \leftarrow \psi - \alpha\,\nabla_\psi\,\mathcal{L}_{\text{post}}$
};

% S4
\node[stepblock, below=0.5cm of s3] (s4) {
    \textbf{S4. Reward Assignment} \hfill \textit{(no parameter updates)}\\[3pt]
    \textbf{for} each simulated $(s, a, s')$ with segment $z$ \textbf{do}\\
    \quad $\tilde{r}_{\text{adv}} \leftarrow f_{\theta,\phi}(s, a, s', z) - \log \pi_\omega(a \mid s, z)$\\
    \quad $\tilde{r}_{\text{MI}} \leftarrow \log q_\psi(z \mid \tau)$ \quad {\footnotesize (standardized)}\\
    \quad $\tilde{r} \leftarrow \tilde{r}_{\text{adv}} + \lambda_{\text{MI}} \cdot \tilde{r}_{\text{MI}}$\\
    \textbf{end for}\\
    {\footnotesize After posterior converges: freeze $q_\psi$, set $\lambda_{\text{MI}} = 0$.}
};

% S5
\node[stepblock, below=0.5cm of s4] (s5) {
    \textbf{S5. Policy Update via PPO} \hfill \textit{(updates $\omega$)}\\[3pt]
    \textbf{for} $z = 1, \ldots, K$ \textbf{do}\\
    \quad $\omega \leftarrow \text{PPO}\!\left(\omega;\; \max_{\pi_\omega} \;\mathbb{E}\!\left[\sum_t \gamma^t\, \tilde{r}\right]\right)$ \quad {\footnotesize (clipped surrogate objective)}\\
    \textbf{end for}
};

% Convergence check
\node[block, below=0.5cm of s5, fill=blue!5] (conv) {
    \textbf{Check:} $\max_{s,a,z} |\pi_\omega^{\text{new}}(a \mid s, z) - \pi_\omega^{\text{old}}(a \mid s, z)| < \varepsilon$ \;? \quad Stop. \quad Else repeat from S1.
};

% Arrows down
\draw[arrow] (init) -- (s1);
\draw[arrow] (s1) -- (s2);
\draw[arrow] (s2) -- (s3);
\draw[arrow] (s3) -- (s4);
\draw[arrow] (s4) -- (s5);
\draw[arrow] (s5) -- (conv);

% Loop arrow
\draw[looparrow, rounded corners=8pt]
    (conv.east) -- ++(1.2,0) |- (s1.east)
    node[pos=0.25, right, font=\small\itshape, text=blue!60] {repeat};

\end{tikzpicture}
\end{center}

\newpage
\section*{At Convergence}

When $D = 1/2$ everywhere (discriminator cannot distinguish observed from simulated):
\[
f_{\theta,\phi}(s, a, s', z) = \log \pi^*(a \mid s, z) = A^*(s, a, z)
\]
where $A^*$ is the advantage function. With the anchors $g_\theta(s, \text{exit}, z) = 0$ for all $s, z$ and $h_\phi(s_{\text{abs}}, z) = 0$, the decomposition is unique:
\[
g_\theta(s, a, z) = r^*(s, a, z), \qquad h_\phi(s, z) = V^*(s, z).
\]

\end{document}
